\documentclass[a4paper]{article}     
\usepackage{amsfonts}
\usepackage{amsmath}
\usepackage{amssymb}
\usepackage{graphicx}
\usepackage{cancel}

\newcommand{\asa}{\Leftrightarrow} %als slechts als
\newcommand{\adR}{\Rightarrow} %als dan Rechts
\newcommand{\adL}{\Leftarrow}
\newcommand{\Lh}{\left(} %Links haakje
\newcommand{\Rh}{\right)}
\newcommand{\Lv}{\left[} %Links vierkanthaakje
\newcommand{\Rv}{\right]}
\newcommand{\La}{\left\{} %Links acolade
\newcommand{\Ra}{\right\}}
\newcommand{\Ras}{\right.} %Rechts acolade stelsel (omdat om stelsels te maken heb je een punt nodig
\newcommand{\pnR}{\rightarrow} %pijl naar Rechts
\newcommand{\limiet}{\lim\limits_}
\newcommand{\schrap}{\cancel}
\newcommand{\schrapb}{\bcancel} %backslash
\newcommand{\schrapk}{\xcancel} %kruis
\newcommand{\vervang}{\cancelto}
\newcommand{\intgrl}{\displaystyle\int} %integraal teken (bij het berekenen van primitieven)

\begin{document}

\noindent \large Auteur: Yoren Collier \\
\noindent \large Studierichting: Informatica\\
\noindent \large Oefeningenreeks 4A nr. 37.14\\

\medskip

\normalsize

$ I =  \intgrl \dfrac{4dx}{x^2 +x +1} dx = \intgrl \dfrac{4dx}{x^2 +x \dfrac{1}{4} + \dfrac{3}{4}} dx = \intgrl \dfrac{4dx}{\Lh x + \dfrac{1}{2} \Rh^2 + \dfrac{3}{4}} dx$\\

Stel: $t = x + \dfrac{1}{2} \adR dt = dx$\\

$= 4\intgrl \dfrac{1}{t^2 + \sqrt{\dfrac{3}{4}}} dt = 4 \cdot \dfrac{1}{\sqrt{\dfrac{3}{4}}} \operatorname{Bgtan} \Lh \dfrac{t}{\sqrt{\dfrac{3}{4}}} \Rh +c$\\

$= \dfrac{8}{\sqrt{3}} \operatorname{Bgtan} \Lh \dfrac{x + \dfrac{1}{2}}{\dfrac{\sqrt{3}}{2}} \Rh +c = \dfrac{8\sqrt{3}}{3} \operatorname{Bgtan} \Lh \dfrac{\Lh 2x + 1 \Rh \sqrt{3}}{3} \Rh +c$\\

\begin{thebibliography}{999}
%\bibitem{a} Referentie
\end{thebibliography}
\end{document} 